\chapter{Language Grammar Design}

\section{The Grammar}

The DSL follows an imperative, statement-oriented structure where a program consists of a sequential list of statements. 
Each statement occupies its own line, with newlines acting as statement terminators rather than semicolons, reducing syntactic noise for novice learners.

The language supports six categories of statements: variable assignments, two forms of iteration, a while loop, a conditional, function calls, and single-line comments. 
Comments are introduced with a double dash (\texttt{--}), a deliberate choice to avoid conflict with the subtraction operator and to remain readable.

Iteration is provided in two forms. 
A count-based loop (\texttt{for N times}) allows repetition without requiring the learner to manage an explicit iterator, which lowers the cognitive barrier for simple repetitive tasks. 
A range-based loop (\texttt{for x from A to B}) introduces the concept of an iterator variable progressively, once the learner is ready for it. 
Both forms use indentation-based block delimitation, consistent with Python-like conventions that children are likely to encounter in subsequent languages.

Expressions follow a standard left-recursive arithmetic structure supporting addition, subtraction, multiplication, and division, with parentheses available for grouping. 
A special dot-notation construct (\texttt{object.property}) is provided to represent Minecraft-specific entities and their attributes, offering a lightweight form of structured data access without introducing full object-oriented semantics.

Identifiers follow conventional alphanumeric rules with underscore support. 
Literals include integers, quoted strings, and boolean values (\texttt{true}/\texttt{false}). 
Whitespace handling is made explicit in the grammar, with optional spacing permitted around operators and delimiters to give learners flexibility in formatting their code.

\subsection{Grammar Notation Conventions}

The grammar is written in Backus--Naur Form (BNF), where non-terminal symbols are enclosed in angle brackets and terminal symbols appear as quoted string literals.
Several special tokens and conventions are used throughout the grammar and are defined as follows:

\begin{itemize}
    \item[\textbf{--}] \textbf{NEWLINE:} Represents a line break character (\texttt{\textbackslash n}), used to terminate statements and signal the end of a logical line.
    \item[\textbf{--}] \textbf{INDENT:} Represents an increase in indentation level at the beginning of a line. The grammar does not prescribe a fixed indentation width; any consistent combination of spaces or tabs is accepted, consistent with learner-friendly formatting conventions.
    \item[\textbf{--}] \textbf{DEDENT:} Represents a decrease in indentation level, signaling the end of an indented block. It is the logical counterpart of \textbf{INDENT}.
    \item[\textbf{--}] \textbf{$\varepsilon$:} Denotes the empty string, used in productions where a symbol may derive nothing, such as optional whitespace or an empty character sequence.
\end{itemize}

The complete formal grammar in BNF notation is provided in Appendix~\ref{appendix:grammar}.

\section{Lexer Implementation}

The lexer is responsible for transforming raw source text into a flat sequence of tokens, which the parser then consumes. It operates line by line, handling both the structural indentation logic and the inline tokenization of each line's content.

\subsection{Token Recognition}

Token recognition is driven by a list of ordered token specifications, each pairing a \texttt{TokenType} with a regular expression pattern. The order is significant --- more specific or longer patterns are listed before shorter ones to prevent incorrect partial matches. For example, the two-character operators \texttt{!=}, \texttt{==}, \texttt{<=}, and \texttt{>=} are specified before their single-character counterparts \texttt{<} and \texttt{>}.

All individual patterns are combined at class initialization into a single master regular expression using alternation, and compiled once into a \texttt{Pattern} object. During tokenization, a \texttt{Matcher} is advanced through the content of each line using \texttt{lookingAt()}, which attempts to match at the current position without requiring a full-string match. The capturing group that produced the match is then used to identify the corresponding \texttt{TokenSpec}.

\subsection{Indentation Handling}

Since the DSL uses indentation-based block delimitation, the lexer is responsible for emitting \texttt{INDENT} and \texttt{DEDENT} tokens explicitly. 
An integer stack tracks the indentation levels encountered so far, initialized with a base level of zero. At the start of each line, the number of leading spaces and tabs is counted. 
If the measured indentation exceeds the top of the stack, an \texttt{INDENT} token is emitted and the new level is pushed. If it is less, \texttt{DEDENT} tokens are emitted and levels are popped until the stack matches the current indentation. 
Blank and whitespace-only lines are skipped entirely and carry no structural meaning.

\subsection{Keyword Resolution and Special Cases}

Identifiers are first matched by a general alphanumeric pattern, then passed through a keyword lookup to determine whether the matched text corresponds to a reserved word such as \texttt{for}, \texttt{while}, \texttt{if}, or \texttt{true}. If a keyword match is found, the corresponding keyword token type is emitted instead of a generic \texttt{IDENT} token.

String literals have their surrounding quotation marks stripped during tokenization, so the parser receives only the inner content. Comments, introduced by a double dash (\texttt{--}), cause the lexer to discard the remainder of the line immediately upon detection. Any character that does not match any known pattern results in an \texttt{ILLEGAL} token carrying the unrecognized character, allowing the parser to produce a meaningful error message rather than failing silently.

Each line concludes with an explicit \texttt{NEWLINE} token. At the end of the source, any remaining open indentation levels are closed with corresponding \texttt{DEDENT} tokens, followed by a final \texttt{EOF} token.

\section{Parser Implementation}

The parser consumes the token sequence produced by the lexer and constructs an Abstract Syntax Tree (AST). 
It is implemented as a hand-written recursive descent parser, where each grammar rule corresponds directly to a parsing method.

The parser maintains a cursor over the token list, advancing through it with \texttt{expect()}, which asserts the current token type and throws a \texttt{ParseException} with the line and column number if the assertion fails.

When the current token is an identifier, a one-token lookahead disambiguates between an assignment and a function call statement. The two forms of \texttt{for} loop are similarly disambiguated by checking whether the identifier after \texttt{for} is followed by the \texttt{from} keyword.

Expression parsing uses a two-level recursive descent encoding operator precedence: \texttt{parseExpression()} handles additive operators, \texttt{parseTerm()} handles multiplicative operators, and \texttt{parseFactor()} handles the base cases. Conditions are parsed separately, supporting negation, comparison operators, and bare identifiers or function calls used as boolean predicates.

All syntax errors are reported as a \texttt{ParseException} carrying the offending token, its literal value, and its position in the source.